     5th INTERNATIONAL MATHEMATICS COMPETITION FOR UNIVERSITY
                                  STUDENTS
                 July 29 - August 3, 1998, Blagoevgrad, Bulgaria

                                                   Second day

                                       PROBLEMS AND SOLUTION


Problem 1. (20 points) Let V be a real vector space, and let f, f1 , f2 , . . . , fk be linear maps from V
to IR. Suppose that f (x) = 0 whenever f1 (x) = f2 (x) = . . . = fk (x) = 0. Prove that f is a linear
combination of f1 , f2 , ..., fk .

Solution. We use induction on k. By passing to a subset, we may assume that f1 , . . . , fk are linearly
independent.
      Since fk is independent of f1 , . . . , fk−1 , by induction there exists a vector ak ∈ V such that f1 (ak ) =
. . . = fk−1 (ak ) = 0 and fk (ak ) 6= 0. After normalising, we may assume that fk (ak ) = 1. The vectors
a1 , . . . , ak−1 are defined similarly to get
                                                           
                                                             1 if i = j
                                                fi (aj ) =
                                                             0 if i 6= j.
                                                                                              Pk
For an arbitrary x ∈ V and 1 ≤ i ≤ k, fi (x−f1 (x)a1 −f2 (x)a2 −· · ·−fk (x)ak ) = fi (x)− j=1 fj (x)fi (aj ) =
fi (x) − fi (x)fi (ai ) = 0, thus f (x − f1 (x)a1 − · · · − fk (x)ak ) = 0. By the linearity of f this implies
f (x) = f1 (x)f (a1 ) + · · · + fk (x)f (ak ), which gives f (x) as a linear combination of f1 (x), . . . , fk (x).

Problem 2. (20 points) Let
                                         3
                                         X                                         1
                      P = {f : f (x) =         ak xk , ak ∈ IR, |f (±1)| ≤ 1, |f (± )| ≤ 1}.
                                                                                   2
                                         k=0

Evaluate
                                                sup max |f 00 (x)|
                                                f ∈P −1≤x≤1

and find all polynomials f ∈ P for which the above “sup” is attained.

Solution. Denote x0 = 1, x1 = − 12 , x2 = 12 , x3 = 1,
                                                          3
                                                          Y
                                                w(x) =            (x − xi ),
                                                          i=0

                                                     w(x)
                                         wk (x) =          , k = 0, . . . , 3,
                                                    x − xk
                                                              wk (x)
                                                  lk (x) =             .
                                                              wk (xk )
Then for every f ∈ P
                                                         3
                                                         X
                                            f 00 (x) =            lk00 (x)f (xk ),
                                                         k=0
                                                              3
                                                              X
                                               |f 00 (x)| ≤          |lk00 (x)|.
                                                              k=0




                                                              1
Since f 00 is a linear function max−1≤x≤1 |f 00 (x)| is attained either at x = −1 or at x = 1. Without loss
of generality let the maximum point is x = 1. Then
                                                                        3
                                                                        X
                                         sup max |f 00 (x)| =                 |lk00 (1)|.
                                         f ∈P −1≤x≤1
                                                                        k=0

In order to have equality for the extremal polynomial f∗ there must hold

                                         f∗ (xk ) = signlk00 (1), k = 0, 1, 2, 3.

It is easy to see that {lk00 (1)}3k=0 alternate in sign, so f∗ (xk ) = (−1)k−1 , k = 0, . . . , 3. Hence f∗ (x) =
T3 (x) = 4x3 − 3x, the Chebyshev polynomial of the first kind, and f∗00 (1) = 24. The other extremal
polynomial, corresponding to x = −1, is −T3 .

Problem 3. (20 points) Let 0 < c < 1 and
                                                     x
                                                     c             for x ∈ [0, c],
                                          f (x) =
                                                     1−x
                                                          1−c       for x ∈ [c, 1].

We say that p is an n-periodic point if
                                                 f (f (. . . f (p))) = p
                                                 | {z }
                                                      n

and n is the smallest number with this property. Prove that for every n ≥ 1 the set of n-periodic points
is non-empty and finite.

Solution. Let fn (x) = f (f (. . . f (x))). It is easy to see that fn (x) is a picewise monotone function and
                       | {z }
                                n
its graph contains 2n linear segments; one endpoint is always on {(x, y) : 0 ≤ x ≤ 1, y = 0}, the other is
on {(x, y) : 0 ≤ x ≤ 1, y = 1}. Thus the graph of the identity function intersects each segment once, so
the number of points for which fn (x) = x is 2n .
    Since for each n-periodic points we have fn (x) = x, the number of n-periodic points is finite.
    A point x is n-periodic if fn (x) = x but fk (x) 6= x for k = 1, . . . , n−1. But as we saw before fk (x) = x
holds only at 2k points, so there are at most 21 + 22 + · · · + 2n−1 = 2n − 2 points x for which fk (x) = x
for at least one k ∈ {1, 2, . . . , n − 1}. Therefore at least two of the 2n points for which fn (x) = x are
n-periodic points.

Problem 4. (20 points) Let An = {1, 2, . . . , n}, where n ≥ 3. Let F be the family of all non-constant
functions f : An → An satisfying the following conditions:
 (1) f (k) ≤ f (k + 1) for k = 1, 2, . . . , n − 1,
 (2) f (k) = f (f (k + 1)) for k = 1, 2, . . . , n − 1.
Find the number of functions in F.

Solution. It is clear that id : An −→ An , given by id(x) = x, does not verify condition (2). Since id is
the only increasing
                     injection on An , F does not contain injections. Let us take any f ∈ F and suppose
that # f −1 (k) ≥ 2. Since f is increasing, there exists i ∈ An such that f (i) = f (i + 1) = k. In view of
(2), f (k) = f (f (i + 1)) = f (i) = k. If {i < k : f (i) < k} = ∅, then taking j = max{i < k : f (i) < k}
                                                                                                         we
get f (j) < f (j + 1) = k = f (f (j + 1)), a contradiction. Hence f (i) = k for i ≤ k. If # f −1 ({l})  ≥ 2
for some l ≥ k, then the similar consideration shows that f (i) = l = k for i ≤ k. Hence # f −1 {i} = 0
or 1 for every i > k. Therefore f (i) ≤ i for i > k. If f (l) = l, then taking j = max{i < l : f (i) < l}
we get f (j) < f (j + 1) = l = f (f (j + 1)), a contradiction. Thus, f (i) ≤ i − 1 for i > k. Let
m = max{i : f (i) = k}. Since f is non-constant m ≤ n − 1. Since k = f (m) = f (f (m + 1)),
f (m + 1) ∈ [k + 1, m]. If f (l) > l − 1 for some l > m + 1, then l − 1 and f (l) belong to f −1 (f (l)) and

                                                                2
this contradicts the facts above. Hence f (i) = i − 1 for i > m + 1. Thus we show that every function f
in F is defined by natural numbers k, l, m, where 1 ≤ k < l = f (m + 1) ≤ m ≤ n − 1.
                                             
                                              k       if i ≤ m
                                     f (i) =   l       if i = m
                                             
                                               i − 1 if i > m + 1.

Then                                                                      
                                                                       n
                                                       #(F) =                  .
                                                                       3


Problem 5. (20 points) Suppose that S is a family of spheres (i.e., surfaces of balls of positive radius)
in IRn , n ≥ 2, such that the intersection of any two contains at most one point. Prove that the set M of
those points that belong to at least two different spheres from S is countable.

Solution. For every x ∈ M choose spheres S, T ∈ S such that S 6= T and x ∈ S ∩ T ; denote by U, V, W
the three components of Rn \ (S ∪ T ), where the notation is such that ∂U = S, ∂V = T and x is the only
point of U ∩ V , and choose points with rational coordinates u ∈ U , v ∈ V , and w ∈ W . We claim that
x is uniquely determined by the triple hu, v, wi; since the set of such triples is countable, this will finish
the proof.
    To prove the claim, suppose, that from some x0 ∈ M we arrived to the same hu, v, wi using spheres
S , T 0 ∈ S and components U 0 , V 0 , W 0 of Rn \ (S 0 ∪ T 0 ). Since S ∩ S 0 contains at most one point and since
  0

U ∩ U 0 6= ∅, we have that U ⊂ U 0 or U 0 ⊂ U ; similarly for V ’s and W ’s. Exchanging the role of x and
x0 and/or of U ’s and V ’s if necessary, there are only two cases to consider: (a) U ⊃ U 0 and V ⊃ V 0 and
(b) U ⊂ U 0 , V ⊃ V 0 and W ⊂ W 0 . In case (a) we recall that U ∩ V contains only x and that x0 ∈ U 0 ∩ V 0 ,
so x = x0 . In case (b) we get from W ⊂ W 0 that U 0 ⊂ U ∪ V ; so since U 0 is open and connected, and
U ∩ V is just one point, we infer that U 0 = U and we are back in the already proved case (a).

Problem 6. (20 points) Let f : (0, 1) → [0, ∞) be a function that is zero except at the distinct points
a1 , a2 , ... . Let bn = f (an ).
                       ∞
                       X
 (a) Prove that if           bn < ∞, then f is differentiable at at least one point x ∈ (0, 1).
                       n=1
                                                                                                         ∞
                                                                                                         P
 (b) Prove that for any sequence of non-negative real numbers (bn )∞
                                                                   n=1 , with                                  bn = ∞, there exists a
                                                                                                         n=1
       sequence (an )∞
                     n=1 such that the function f defined as above is nowhere differentiable.

Solution
                                                                                ∞
                                                                                P
   a) We first construct a sequence cn of positive numbers such that cn → ∞ and   cn bn < 21 . Let
                                                                                                                  n=1
       ∞
       P
B=         bn , and for each k = 0, 1, . . . denote by Nk the first positive integer for which
     n=1

                                                         ∞
                                                         X             B
                                                               bn ≤       .
                                                                       4k
                                                        n=Nk

                  k
             2
Now set cn = 5B for each n, Nk ≤ n < Nk+1 . Then we have cn → ∞ and
                      ∞               ∞                          ∞    ∞                   ∞
                      X               X      X                   X 2k X                   X 2k           B   2
                            cn bn =                    cn bn ≤                     bn ≤              ·      = .
                      n=1
                                                                       5B                       5B       4k  5
                                      k=0 Nk ≤n<Nk+1             k=0        n=Nk          k=0
                                                                                        P
   Consider the intervals In = (an − cn bn , an + cn bn ). The sum of their lengths is 2 cn bn < 1, thus
there exists a point x0 ∈ (0, 1) which is not contained in any In . We show that f is differentiable at x0 ,


                                                                 3
and f 0 (x0 ) = 0. Since x0 is outside of the intervals In , x0 6= an for any n and f (x0 ) = 0. For arbitrary
x ∈ (0, 1) \ {x0 }, if x = an for some n, then

                                  f (x) − f (x0 )   f (an ) − 0    bn     1
                                                  =             ≤       = ,
                                      x − x0        |an − x0 |    cn bn  cn

otherwise f (x)−f
              x−x0
                   (x0 )
                         = 0. Since cn → ∞, this implies that for arbitrary ε > 0 there are only finitely many
x ∈ (0, 1) \ {x0 } for which
                                               f (x) − f (x0 )
                                                               <ε
                                                   x − x0
does not hold, and we are done.
    Remark. The variation of f is finite, which implies that f is differentiable almost everywhere .
    b) We remove the zero elements from sequence bn . Since f (x) = 0 except for a countable subset of
(0, 1), if f is differentiable at some point x0 , then f (x0 ) and f 0 (x0 ) must be P
                                                                                     0.
    It is easy to construct a sequence βn satisfying 0 < βn ≤ bn , bn → 0 and ∞         n=1 βn = ∞.
    Choose the numbers a1 , a2 , . . . such that the intervals In = (an − βn , an + βn ) (nP= 1, 2, . . .) cover
each point of (0, 1) infinitely many times (it is possible since the sum of lengths is 2 bn = ∞). Then
for arbitrary x0 ∈ (0, 1), f (x0 ) = 0 and ε > 0 there is an n for which βn < ε and x0 ∈ In which implies

                                         |f (an ) − f (x0 )|   bn
                                                             >    ≥ 1.
                                             |an − x0 |        βn




                                                       4
